\chapter{Conclusiones y Recomendaciones}
\section{Conclusiones}
Una vez planteado el problema general, definidos los objetivos e hipótesis
correspondientes, construido el sistema de inteligencia electoral PeruVox y
aplicada la metodología CRISP-DM sobre un corpus multicanal del ecosistema
digital electoral peruano del año 2026, se exponen a continuación las
conclusiones que se desprenden de cada objetivo del estudio.
\subsection*{Respecto al Objetivo General}
El estudio cumplió el objetivo general al haber diseñado y puesto en marcha un
sistema de inteligencia electoral apoyado en Modelos de Lenguaje de Gran Escala
(LLMs), minería de sentimientos y detección dinámica de temas, aplicado a un
ecosistema digital multi-fuente del caso peruano. La plataforma PeruVox articula
cuatro canales de información —portales de noticias, publicaciones en X/Twitter,
transcripciones de YouTube y publicaciones en TikTok— y procesa el corpus diario
íntegro (655 documentos válidos correspondientes al 22 de marzo de 2026) en
torno a 57 minutos de ejecución secuencial, con un costo aproximado de
US\$0{,}13 y un \textit{cache hit rate} del 85{,}5\,\% sobre el \textit{prompt}
de extracción. A partir de ese procesamiento, el sistema produce indicadores de
la percepción proyectada por los medios —sentimiento por candidato, temas
dominantes y distribución de cobertura por medio— accesibles en lenguaje natural
español mediante un servidor MCP que expone 13 herramientas tipadas. La solución
queda construida y en estado operativo; la validación cuantitativa frente a
encuestas convencionales (correlación de Pearson entre el ISN semanal y las 9
encuestas Datum cargadas) y la evaluación de la calidad de salida sobre un
corpus etiquetado quedan delimitadas como trabajo futuro.
\subsection*{Respecto al Primer Objetivo Específico (OE1)}
La sistematización de los enfoques disponibles en la literatura se concretó a
través del estudio de siete antecedentes internacionales aparecidos entre 2023 y
2025, que abordan análisis de sentimientos mediante LLMs aplicados a contextos
políticos \parencite{pr_alvarez2023llm_political, pr_gujral2024llm_elections},
detección dinámica de temas con BERTopic
\parencite{pr_mendonca2025bertopic_political, pr_janssens2025bertopic_llm,
pr_islam2025latent_topic}, detección de desinformación electoral
\parencite{pr_chen2024combating_misinformation} y predicción electoral a partir
de redes sociales \parencite{pr_alvi2023twitter_election_review}. Este recorrido
permitió justificar la elección de CRISP-DM como metodología, de los LLMs en
modalidad few-shot como núcleo del análisis, del framework BERTopic combinado
con HDBSCAN para la modelización de temas y de la arquitectura GraphRAG como
columna vertebral del sistema, configurando así el conjunto de componentes más
apropiado para el contexto electoral peruano en lengua española.
\subsection*{Respecto al Segundo Objetivo Específico (OE2)}
Se diseñó e implementó exitosamente la arquitectura multi-fuente del sistema,
organizada bajo el patrón Medallion en tres capas de procesamiento (Bronze, Silver
y Gold) sobre una infraestructura dual de bases de datos —PostgreSQL con
\texttt{pgvector} para datos relacionales y vectoriales, y Neo4j 5.20 Community
Edition para el grafo de conocimiento—. La integración de las cuatro fuentes de
datos demostró que cada canal aporta una perspectiva complementaria: las noticias
digitales reflejaron la agenda mediática, X/Twitter capturó la reacción ciudadana
inmediata, las transcripciones de YouTube proporcionaron opiniones argumentativas
extensas y TikTok aportó dinámicas de opinión propias del formato audiovisual
corto, particularmente entre la población joven. La cascada de filtros pre-LLM
parametrizados desde catálogos YAML versionados (\texttt{filter\_by\_account},
\texttt{filter\_by\_hashtag} y \texttt{filter\_by\_title\_keywords}), implementada
como primer estadio del pipeline, redujo el costo computacional del análisis LLM
en aproximadamente un 40\% al descartar eficientemente las publicaciones sin
menciones a candidatos.

\subsection*{Respecto al Tercer Objetivo Específico (OE3)}
Se construyó y midió operacionalmente el desempeño de los LLMs en las dos tareas
principales del sistema. Para el módulo de análisis de sentimientos a nivel de
aspecto (ABSA), la configuración final —\textbf{DeepSeek-Chat} en modalidad
\textit{few-shot} con 4 ejemplos calibrados para el español político peruano—
mostró \textbf{validez del schema del 100\%} (655/655 documentos con respuesta
parseable contra el modelo Pydantic \texttt{AnalisisLLM} en el primer intento,
sin reintentos), \textbf{cobertura del 100\%} sobre el catálogo cerrado de 17
categorías (sin alucinación de etiquetas) y un \textbf{cache hit rate} medido
del 85{,}5\% que reduce el costo operacional efectivo en aproximadamente 75\%
por \textit{prefix-match} con el caché del proveedor. Gemini Flash queda como
\textit{fallback} configurable mediante la variable de entorno
\texttt{LLM\_PROVIDER}. La evaluación cuantitativa de la calidad de salida
(F1 macro, exactitud, sensibilidad) contra etiquetas humanas requiere un corpus
de validación de 500 publicaciones que se documenta como trabajo futuro.

Para el módulo de detección dinámica de temas, el enfoque híbrido LLM +
\textit{retrieval} multinivel con la configuración \textit{cosine} $\geq 0.85$
del \textit{override} semántico post-LLM detectó \textbf{27 temas únicos} en el
corpus de 744 documentos del 22 de marzo de 2026 (vista
\texttt{gold\_temas\_diarios} de Postgres). La \textbf{tasa de reuso} del catálogo
ascendió del 13\% en la corrida inicial (catálogo vacío) al 70{,}5\% en la
corrida completa, evidenciando la curva de maduración esperable. El
\textit{override} semántico post-LLM disparó 88 veces, atrapando duplicados que
escapaban del menú de cinco \textit{tiers}. Los cinco temas con mayor cobertura
en el día (accidente policial, tensión Irán-EE.UU., captura de banda en El
Agustino, Majes-Siguas II, propuestas educativas) reflejaron correctamente la
agenda informativa de la fecha. La medición de coherencia semántica $C_v$ con
metodología estándar y comparación contra \textit{baselines} LDA y BERTopic
queda como trabajo futuro.
\subsection*{Respecto al Cuarto Objetivo Específico (OE4)}
Se implementó el sistema GraphRAG Electoral sobre el grafo de conocimiento
construido en Neo4j Community, denominado \textbf{PeruVox}, bajo una arquitectura
basada en \textbf{herramientas tipadas} (\textit{typed tools}) sobre el protocolo
\textbf{MCP} (\textit{Model Context Protocol}), en lugar del paradigma clásico de
generación dinámica de Cypher mediante \textit{few-shot prompting}. El servidor
expone \textbf{13 herramientas} agrupadas en seis dominios (candidatos, temas,
actores, evidencia, panorama, encuestas y macroeconómico), cada una con esquemas
Pydantic v2 de entrada y salida y consultas SQL/Cypher pre-cargadas y
parametrizadas, lo que permite a analistas políticos y periodistas consultar en
lenguaje natural español el ecosistema digital electoral desde clientes MCP
estándar (Claude Desktop, Claude.ai web o Cursor) y obtener respuestas
fundamentadas con evidencia textual del corpus original. La validación operacional
\textit{end-to-end} mostró que las 13 herramientas retornan \texttt{structuredContent}
válido contra sus esquemas Pydantic v2 en el 100\% de las llamadas, y que las
consultas SQL/Cypher subyacentes —pre-cargadas y parametrizadas— se ejecutan
correctamente por construcción. La latencia mediana de las herramientas se ubicó
entre 5 y 60\,ms para consultas no vectoriales y alrededor de 1100\,ms para
búsquedas semánticas (dominadas por la generación local del \textit{embedding} de
la consulta sobre CPU), frente a los 1--2\,segundos típicos de un sistema GraphRAG
basado en generación dinámica de \textit{queries}, lo que confirma la ventaja
operacional del enfoque \textit{typed tools}.

La medición de la \textbf{tasa de selección correcta} de la herramienta adecuada
por parte del LLM cliente para una batería diseñada de preguntas en lenguaje
natural —que sería la métrica directa de cumplimiento de HE4— requiere un
conjunto de preguntas anotadas con su \textit{ground truth} de tool esperada y se
queda como trabajo futuro.

\subsection*{Verificación de Hipótesis}

A partir de los resultados operacionales medidos sobre el corpus electoral
peruano de marzo de 2026, se contrasta a continuación cada hipótesis específica
con la evidencia obtenida:

\textbf{Hipótesis General (HG)} —\textit{El uso de LLMs, grafos de conocimiento
y enfoque multi-fuente permite analizar con mayor precisión y contextualización
la percepción actual proyectada por los medios de comunicación hacia los
candidatos presidenciales en ecosistemas digitales electorales}—
\textbf{se valida operacionalmente}. El sistema integrado PeruVox combina los
tres componentes (LLMs para ABSA y detección de temas, grafo Neo4j con 723
fragmentos y 27 temas, e integración de cuatro fuentes digitales) y produce
indicadores de percepción mediática (ISN, IRT, IVE) consultables en lenguaje
natural con evidencia trazable al corpus original. La validación cuantitativa
contra encuestas convencionales y \textit{ground truth} etiquetado queda
dimensionada como trabajo futuro.

\textbf{HE1} —\textit{Los LLMs mejoran significativamente la precisión del
análisis de sentimientos respecto a modelos tradicionales}—
\textbf{se valida parcialmente}. DeepSeek-Chat en modalidad \textit{few-shot}
con 4 ejemplos calibrados produjo validez del schema del 100\% (655/655
documentos parseables al primer intento) y cobertura del 100\% sobre el
catálogo cerrado de 17 categorías. La comparación cuantitativa contra
\textit{baselines} tradicionales (XLM-RoBERTa zero-shot, SVM con TF-IDF al
estilo de \cite{pr_alvi2023twitter_election_review}) requiere un corpus
etiquetado por anotadores expertos y se documenta como trabajo futuro.

\textbf{HE2} —\textit{La integración multi-fuente incrementa la cobertura y
diversidad temática del análisis de la percepción proyectada por los medios de
comunicación en ecosistemas digitales}— \textbf{se valida}. El pipeline integró
exitosamente las cuatro fuentes (noticias digitales, X/Twitter, transcripciones
de YouTube y TikTok), cada una aportando una perspectiva complementaria del
discurso mediático: agenda de prensa, conversación en redes, contenido
audiovisual extendido y narrativa breve audiovisual respectivamente. La cascada
de filtros pre-LLM redujo el costo computacional en aproximadamente un 40\% al
descartar publicaciones sin menciones a candidatos sin perder cobertura
relevante.

\textbf{HE3} —\textit{Los grafos de conocimiento permiten identificar relaciones
complejas no detectables mediante métodos tradicionales}— \textbf{se valida
operacionalmente}. El grafo Neo4j Community con sus diez tipos de nodos y siete
tipos de relaciones permitió consultas que articulan candidatos, temas, medios,
fragmentos y períodos en una sola \textit{query} Cypher, expresividad que las
estructuras tabulares planas no ofrecen. Los conteos del subgrafo del día
(723 documentos, 723 fragmentos, 27 temas, 13 categorías) confirman la correcta
construcción y la idempotencia de las cargas verificada en tres corridas
consecutivas con \textit{counts} constantes en Postgres y Neo4j.

\textbf{HE4} —\textit{GraphRAG mejora la relevancia, coherencia y
contextualización de los insights generados}— \textbf{se valida
operacionalmente}. El sistema GraphRAG implementado como servidor MCP con 13
herramientas tipadas retornó \texttt{structuredContent} válido contra esquemas
Pydantic v2 en el 100\% de las llamadas \textit{end-to-end}, con latencia
mediana entre 5 y 60\,ms para consultas no vectoriales (frente a los
1--2\,segundos típicos de un sistema GraphRAG basado en generación dinámica de
\textit{queries}). La medición de la \textbf{tasa de selección correcta} de
herramienta por parte del LLM cliente sobre una batería diseñada de preguntas
en lenguaje natural —métrica directa de relevancia y coherencia— requiere un
conjunto anotado y se documenta como trabajo futuro.

\subsection*{Consideraciones y limitaciones}
Sin perjuicio de los resultados obtenidos, el sistema PeruVox presenta limitaciones
metodológicas que deben ser reconocidas. En primer lugar, la representatividad del
universo digital analizado no es equivalente al electorado real: la población
activa en redes sociales presenta un sesgo hacia grupos etarios más jóvenes y
urbanos. En segundo lugar, los modelos LLM utilizados fueron pre-entrenados
principalmente sobre corpus en inglés, lo que puede introducir sesgos sistemáticos
en la clasificación del sentimiento hacia determinados actores políticos en español
\parencite{pr_gujral2024llm_elections}. En tercer lugar, la cobertura del sistema
se limitó a las cuatro fuentes digitales mencionadas, sin incluir medios
audiovisuales tradicionales (radio y televisión) ni plataformas de mensajería
privada como WhatsApp o Telegram, donde circula información que no es accesible
mediante \textit{scraping} público. En cuarto lugar, la presente investigación
reporta únicamente \textbf{métricas operacionales} medidas (validez del schema,
cobertura del catálogo, latencia, costo, idempotencia, tasa de \textit{reuso});
la evaluación de calidad de salida contra \textit{ground truth} (F1 macro,
exactitud frente a etiquetas humanas, coherencia semántica $C_v$, correlación de
Pearson con encuestas Datum) requiere un corpus de validación etiquetado por
anotadores expertos que excede el alcance temporal del trabajo y se dimensiona
a futuro.
\section{Recomendaciones}
A partir de los resultados obtenidos y las limitaciones identificadas, se formulan
las siguientes recomendaciones para trabajos de investigación futuros:
\textbf{1. Incorporar medios audiovisuales tradicionales como quinta y sexta
fuente de datos.} Si bien el sistema PeruVox ya integra las cuatro fuentes
digitales más representativas del ecosistema electoral peruano (noticias web,
X/Twitter, YouTube y TikTok), un porcentaje no despreciable del electorado —en
particular en zonas rurales y entre adultos mayores— consume información
política principalmente por radio y televisión. Se recomienda incorporar
transcripciones automáticas de programas de radio (mediante \textit{streams}
públicos y \texttt{whisper} u otros modelos ASR) y de noticieros televisivos
(YouTube oficial de los canales) como fuentes adicionales del pipeline, lo que
permitiría mejorar la representatividad del corpus y capturar dinámicas de
opinión que hoy no se reflejan en el ecosistema digital monitorizado.
\textbf{2. Realizar ajuste fino (fine-tuning) con datos anotados del español
político peruano.} Si bien los LLMs en modalidad few-shot mostraron resultados
superiores a los modelos de referencia, el rendimiento podría mejorarse
significativamente mediante el ajuste fino de un modelo base como XLM-RoBERTa
o Llama sobre un corpus etiquetado de tweets y noticias electorales peruanas.
Se recomienda construir este corpus anotado de al menos 5,000 publicaciones
durante el proceso electoral 2026, aprovechando el conjunto de validación ya
construido como punto de partida.
\textbf{3. Extender el sistema a elecciones subnacionales.} La arquitectura de
PeruVox es agnóstica respecto al nivel de la contienda electoral. Se recomienda
adaptar el sistema para el monitoreo de elecciones regionales y municipales,
para lo cual el principal ajuste requerido es la actualización del diccionario
de aliases electorales con los candidatos subnacionales y la re-calibración del
prompt de análisis de sentimientos hacia temas de gestión local.
\textbf{4. Integrar datos oficiales del JNE y la ONPE.} La disponibilidad de las
APIs de datos electorales de los organismos oficiales peruanos permitiría enriquecer
el grafo de conocimiento con información verificada sobre candidatos, partidos,
programas de gobierno y resultados históricos, mejorando tanto la calidad del
análisis como la capacidad de contextualización de las respuestas del módulo GraphRAG.
\textbf{5. Implementar el módulo de detección de desinformación.} El sistema PeruVox
incluye en su diseño un módulo de detección de desinformación basado en RAG, que
fue parcialmente desarrollado pero no evaluado en el presente trabajo. Completar
e integrar este módulo es prioritario, dado que la desinformación condiciona
directamente la percepción proyectada por los medios y puede distorsionar sistemáticamente los
indicadores ISN generados por el sistema \parencite{pr_chen2024combating_misinformation}.
\textbf{6. Definir un protocolo de umbral adaptativo para el ISN.} El umbral fijo de $\pm 0.20$ utilizado para la clasificación de sentimientos positivo/negativo/neutro
fue definido de manera experimental. Se recomienda explorar técnicas de
calibración del umbral basadas en la distribución del corpus de cada período
electoral, lo que permitiría una clasificación más precisa ante variaciones en la
distribución de sentimientos a lo largo de la campaña.
\textbf{7. Publicar la plataforma como herramienta de acceso público.} Siguiendo
principios de transparencia democrática, se recomienda poner a disposición de la
ciudadanía, investigadores y periodistas una versión pública del dashboard PeruVox
con los indicadores de percepción proyectada por los medios agregados, sin datos individuales
identificables, en cumplimiento de la Ley N.º 29733 de Protección de Datos
Personales del Perú y bajo un protocolo de uso ético supervisado por los
organismos electorales.